\documentclass{article}
\usepackage[T1]{fontenc}
\usepackage{lmodern}
\usepackage{graphicx}
\usepackage{amsmath,amssymb}
\usepackage[a4paper,margin=2cm]{geometry}
\usepackage[absolute,overlay]{textpos}
\setlength{\TPHorizModule}{1mm}
\setlength{\TPVertModule}{1mm}
\usepackage[indonesian]{babel}
\usepackage{hyperref}
\setlength{\emergencystretch}{2em}
% unit: Halaman 18

\begin{document}

% === Halaman 18 (llm) ===

Contoh di dalam Vigenere Cipher:

\begin{itemize}
    \item Enkripsi setiap huruf plainteks:
    \[ C = (P + K) \pmod{26} \]
    \item Penyerang adaptif:
    \begin{itemize}
        \item Pertama masukkan plainteks = AAAA (=0000), maka ciphertext C = K.
        \item Setelah tahu K, ia coba plainteks lain untuk memverifikasi, misalnya BBBB (=1111), hasilnya menegaskan hasil enkripsi dengan K yang sama.
        \item Dengan informasi itu, cipherteks lain bisa langsung dipecahkan.
    \end{itemize}
\end{itemize}

\end{document}
